% Generated by figure/tools/make_provenance.py -- edit the script, not this file.
\begingroup\footnotesize\sloppy
\setlength{\LTcapwidth}{\textwidth}
\begin{longtable}{@{}>{\raggedright\arraybackslash}p{0.05\textwidth}>{\raggedright\arraybackslash}p{0.2\textwidth}>{\centering\arraybackslash}p{0.05\textwidth}>{\raggedright\arraybackslash}p{0.23\textwidth}>{\raggedright\arraybackslash}p{0.16\textwidth}>{\raggedright\arraybackslash}p{0.16\textwidth}@{}}
\caption{Full provenance of every constraint and stated result, with the corrections made while tracing them. Basis: D, definition; C, convention; M, derived and checked symbolically (Table~\ref{tab:symbolic-checks}); Q, quoted; A, adapted. Evidence: full text (with location), abstract, title, or not accessed. Where the parameter table of the model builder uses other symbols than the main text, the rows follow the builder (for example $\mu$ for $n$ in T1.3, $\nu$ for $\mu$ in T2.10, $w$ for $\lambda$ in T2.14, $\psi$ for $\theta$ in T3.3, $\eta$ and $C$ for $n$ and $k$ in T3.5--T3.6, $\zeta$ for $\eta$ in $\Phi_y$.2). Last update: 19 October 2026.}\label{tab:provenance-full}\\
\toprule ID & Statement & Basis & Explanation and source & Checks and evidence & Correction history \\ \midrule \endfirsthead
\toprule ID & Statement & Basis & Explanation and source & Checks and evidence & Correction history \\ \midrule \endhead
\bottomrule \endfoot
\multicolumn{6}{@{}l}{\textit{Prey growth ($T_1$)}}\\[2pt]
G1 & $r>0$ & D & Intrinsic growth rate: the prey increases at low density without predators. & -- & -- \\
G2 & $K>0$ & D & Carrying capacity. & -- & -- \\
G3 & $\theta>0$; $\theta=1$ is the logistic law & D, M & Shape exponent of density dependence \cite{gilpin1973}. & \checkmark\,S1 & -- \\
G4 & $\tau^{x}\ge 0$ & D & Delay in self-limitation \cite{hutchinson1948}. & -- & -- \\
G5 & Without predators, $K$ is locally asymptotically stable iff $r\tau^{x}<\pi/2$; inconclusive at equality & M & Roots of $\lambda+re^{-\lambda\tau}=0$: $-r$ at $\tau=0$; imaginary roots only at $\pm ir$ with $\sin r\tau=1$; transversal crossing at $r\tau=\pi/2$. & \checkmark\,S6 & ``Stable iff $r\tau<\pi/2$\textquotedblright{} did not separate local from global stability; now split into G5 and G5b. Second evaluation: the builder states G5 and G5b only when no harvest, immigration, seasonal forcing, prey competition, diffusion or noise changes the prey-only equation. \\
G5b & Without predators, $K$ attracts all positive solutions for $r\tau^{x}\le\pi/2$ & A & $x=K[1+z(t/\tau)]$ gives Wright's equation with $\alpha=r\tau$ (derived); Wright's conjecture proved in \cite{vandenberg2018}. & \checkmark\,S6b; \cite{vandenberg2018}: full text (arXiv:\allowbreak 1704.00029v1, Thm.~1.5) & Added at review; restricted as G5 at the second evaluation. \\
G6 & Gompertz: $r\ln(K/x)\to\infty$ as $x\to0^{+}$; not Lipschitz at 0 but solutions unique; $x=0$ neither reached nor left in finite time unless immigration, constant yield or herd behaviour act at $x=0$ & M & Osgood: $\int_0 dx/[rx\ln(K/x)]=\infty$ \cite{gompertz1825}. Immigration gives $\dot x=I>0$ at $x=0$; constant yield and the herd term $x^{p}$ reach 0 in finite time. & \checkmark\,S7, S7b, S8d, S21 & Uniqueness statement added at review; exceptions added at the second evaluation. \\
G7 & Strong Allee: $0<A<K$ & M & Per-capita growth $r(1-x/K)(x-A)/K$ has value $-rA/K$ at 0 and zeros $A,K$ \cite{allee1931}. & \checkmark\,S2a, S2c & -- \\
G8 & Weak Allee: $-K<A\le 0$; $A\le -K$ removes the effect & M, C & Slope $r(K+A)/K^{2}$ at 0. The borderline $A=0$ is grouped with the weak case by convention. & \checkmark\,S2b, S2d & -- \\
G9 & Strong Allee, prey alone: origin locally stable, outcome depends on initial density; with prey immigration a small stable equilibrium replaces the origin & M & Follows from $g(0)<0$; with immigration $f(0)=I>0$ and the implicit function theorem. & \checkmark\,S2a, S32a & Restricted to ``without predators\textquotedblright{}; the immigration case added at the second evaluation. \\
G9b & Strong Allee with a predator: heteroclinic loops and a threshold curve separating overexploitation from coexistence & Q & \cite{wang2011allee}. & \cite{wang2011allee}: abstract & -- \\
G10 & Additive Allee: strong iff $b<K$ and $rb<\mu<r(K+b)^{2}/(4K)$; weak iff $rb^{2}/K<\mu\le rb$ (the borderline $\mu=rb$ grouped with the weak case) & M & $g(0)=r-\mu/b$, $g'(0)=-r/K+\mu/b^{2}$, $\max_{x\ge0}r(1-x/K)(x+b)$; additive form from \cite{dennis1989}. & \checkmark\,S3a, S3b, S3c, S3d, S3e & ``Strong iff $\mu>rb$\textquotedblright{} was wrong: $r=K=b=1$, $\mu=2$ gives decline at every density (S3e). \\
G11 & Fear: $r>d$ & M & Low-density growth rate of the prey alone is $r-d$; model of \cite{wang2016fear}. & \checkmark\,S4 & -- \\
G12 & $d\ge 0$, $c^{x}>0$, $k\ge 0$ ($k=0$: no fear) & D & Death rate, self-limitation and fear level \cite{wang2016fear}. & -- & -- \\
G13 & Perceived predation risk alone cut offspring production by 40\% & Q & Field experiment on song sparrows \cite{zanette2011}. & \cite{zanette2011}: abstract & -- \\
G13b & Fear model: high fear excludes periodic solutions; low fear allows subcritical Hopf bifurcation and bistability & Q & \cite{wang2016fear}. & \cite{wang2016fear}: abstract & -- \\
G14 & $0\le\varepsilon<1$; $\omega=2\pi/T>0$ & M & $\min_{t} r(1+\varepsilon\sin\omega t)=r(1-\varepsilon)$. & \checkmark\,S5 & -- \\
G15 & Seasonal forcing: periodic solutions, multiple attractors and chaos & D, Q & Non-autonomous by construction; periodic solutions \cite{cushing1977}; multiple attractors and chaos \cite{rinaldi1993}. & \cite{cushing1977}: title; \cite{rinaldi1993}: title & ``Quasi-periodic\textquotedblright{} removed: not supported by the titles. \\
G16 & $\alpha_{ii'}\ge 0$ & D & Interspecific competition between prey. & -- & -- \\
G17 & Toxicant (T1.10): the prey alone persists only if $r_0-r_1C_0>0$ & M & Prey alone: $x=x_0\exp[(r_0-r_1C_0)t]$, which tends to $0$ when $r_0-r_1C_0<0$; model of \cite{hallam1983}. & \checkmark\,S42c & Added 19 October 2026 (audit of the entries table). \\
\addlinespace
\multicolumn{6}{@{}l}{\textit{Predation ($T_2$)}}\\[2pt]
P1 & $a>0$ & D & Attack rate (search efficiency) \cite{holling1959}. & -- & -- \\
P2 & $h>0$; maximal intake $1/h$ & M & $\lim_{x\to\infty} ax/(1+ahx)=1/h$ \cite{holling1959}. & \checkmark\,S9 & -- \\
P3 & One handling time shared by all food items: $D_j=1+\sum_k a_kh_kx_k$ & M & Time budget $T=T_s(1+\sum_k a_kh_kx_k)$; applies to prey, predators eaten and additional food alike. & \checkmark\,S12 & -- \\
P4 & Holling IV peaks at $\sqrt{G}$, $\sqrt{G}/(1-\phi)$ or $\xi+\sqrt{G}$ & M & Group defence \cite{andrews1968,freedman1986}. & \checkmark\,S10a, S10b, S10c & The builder stated $\sqrt{G}$ for every case; corrected for the two refuge types. \\
P5 & Holling IV, one prey and one predator whose per-capita growth depends on prey density only: at most two interior equilibria; Bogdanov--Takens bifurcation & M, Q & Equilibria solve a quadratic; Bogdanov--Takens \cite{xiao2001siam}. & \checkmark\,S10d; \cite{xiao2001siam}: abstract & Source added. Second evaluation: the builder stated it for every model with Holling IV, food webs included; now only under these conditions. \\
P6 & Ivlev: $\rho,\lambda>0$ & D & \cite{ivlev1961}. & -- & -- \\
P7 & Beddington--DeAngelis, Crowley--Martin: $c\ge 0$, $c=0$ is Holling II & M & \cite{beddington1975,deangelis1975,crowley1989}. & \checkmark\,S11a & -- \\
P8 & Hassell--Varley: $0<\gamma<1$ & M, C & $\gamma=0$: Holling II; $\gamma=1$: ratio dependence; the interval is a convention \cite{hassell1969}. & \checkmark\,S11b & -- \\
P9 & Mutual interference: $0\le\nu<1$; for $\nu>0$ solutions from $y=0$ are not unique, unless a predator Allee effect, a gestation delay, predator immigration or a Leslie--Gower predator equation keeps them unique & M & $\nu<1$ makes predation vanish without predators \cite{hassell1971}; explicit second solution from $y=0$. Predator Allee effect: gain of order $y^{2-\nu}$; delay: the gain uses $y(t-\tau)$ (method of steps); immigration: $\dot y>0$ at $y=0$; Leslie--Gower: logistic in $y$. & \checkmark\,S11c, S11d, S33 & Non-uniqueness added at review; exceptions added at the second evaluation. \\
P10 & Ratio dependence: Lipschitz after extension by 0, not differentiable at $(0,0)$; in the basic model with one logistic prey and one predator the origin can attract all orbits & M, Q & Bounded partial derivatives; direction-dependent derivative. Sectors at the origin \cite{xiao2001rd}; global stability of the origin \cite{kuang1998}. & \checkmark\,S28, S28b; \cite{xiao2001rd}: abstract; \cite{kuang1998}: abstract & ``Undefined at the origin\textquotedblright{} sharpened; ``blow-up techniques\textquotedblright{} removed (not in the abstracts). Second evaluation: the results on the origin are now stated only for the basic model. \\
P11 & Herd behaviour: $0<p<1$; $p=1/2$ for a planar herd & M & Edge of a disc-shaped herd scales as area$^{1/2}$ \cite{ajraldi2011}. & \checkmark\,S8c; \cite{ajraldi2011}: abstract & -- \\
P12 & Herd: not Lipschitz at 0; unique forward in time; prey can vanish in finite time, so not unique backward (not with prey immigration; with a fixed-number refuge the singular point is $x=\xi$) & M & $\dot x\le rx$; $\frac{d}{dt}x^{1-p}\to-(1-p)ay<0$; immigration gives $\dot x=I>0$ at $x=0$. & \checkmark\,S8a, S8b, S8d, S30 & ``Uniqueness fails\textquotedblright{} was imprecise: forward uniqueness holds. Exceptions added at the second evaluation. \\
P13 & $0\le\phi<1$; $\xi\ge 0$ & D & Protected fraction or number of prey \cite{maynardsmith1974,gonzalez2003}. & -- & -- \\
P14 & Fixed-number refuge ($\xi>0$) stabilises the Lotka--Volterra centre, a proportional one does not; refuges can destabilise other models & M, A & Proportional refuge rescales $a$; fixed number gives trace $-r\xi/(x^{*}-\xi)$. Attributed to \cite{maynardsmith1974}; caveat \cite{mcnair1986}. & \checkmark\,S29a, S29b; \cite{maynardsmith1974}: not accessed; \cite{mcnair1986}: abstract & Earlier note replaced; McNair added. \\
P15 & $0\le\kappa<1$ & D & Search efficiency lost to habitat complexity \cite{bairagi2011}. & -- & -- \\
P16 & $w\ge 0$; for $w>0$ with one food item, cooperation gives a component Allee effect for Holling I--IV and herd responses; with Beddington--DeAngelis iff $w>ac$; with Crowley--Martin near $y=0$ iff $w>ac(1+ahx)$, never at high $y$; never with ratio dependence; with Hassell--Varley or mutual interference only above a predator density; with additional food iff $e(1+a^{F}h^{F}\Lambda)>e^{F}a^{F}\Lambda h$ & M & Sign of the $y$-derivative of the per-capita gain. With Leslie--Gower predators the gain does not contain the response. & \checkmark\,S17, S31a, S31b, S31c, S31d, S31e, S31f & Second evaluation: the note claimed a component Allee effect for every response; interference can outweigh cooperation. \\
P17 & Cooperation can produce demographic Allee effects and predator collapse (stated where the per-capita gain increases with predator density) & Q & \cite{teixeiraalves2017}. & \cite{teixeiraalves2017}: abstract & ``Bistability\textquotedblright{} replaced by the abstract's wording; restricted as P16 at the second evaluation. \\
P18 & Additional food: loss $ax/(1+ahx+a_Ah_AA)$, gain $(eax+e_Aa_AA)/(1+ahx+a_Ah_AA)$; $\eta=a_A/a$, $\alpha=h_A/h$, $e_A=e$ in Srinivasu et al. & M, A & Time budget (derived); quality--quantity parameterisation and the role of handling times from \cite{srinivasu2007}. & \checkmark\,S12, S12b, S12c; \cite{srinivasu2007}: abstract & The form with one handling time ($\alpha=1$) and $e_A=e$ was replaced by the general form. \\
P19 & Holling II (T2.2): a predator with death rate $m$ persists only if $e/h>m$ & M & Per-capita gain $e\,ax/(1+ahx)-m<e/h-m$ for all $x$, so the predator declines at every prey density when $e/h\le m$ (intake saturates at $1/h$, S9). & \checkmark\,S9, S42f & Added 19 October 2026 (audit of the entries table). \\
P20 & Power-law interactions (T2.16), $\beta<1$: the prey loss $axy^{\beta}$ is not Lipschitz at $y=0$, but the predator equation keeps the factor $y$, so $y=0$ is invariant and solutions from $y=0$ are unique (unlike T2.10) & M & $\dot y=y(ex^{p}-m-\delta y^{q})$ vanishes at $y=0$ with a finite $y$-derivative there; Gronwall then forces $y\equiv0$. & \checkmark\,S43 & Added 20 October 2026 (referee report: an earlier draft grouped T2.16 with the non-unique case). \\
\addlinespace
\multicolumn{6}{@{}l}{\textit{Predator ($T_3$, $T_4$)}}\\[2pt]
N1 & $e>0$; $e<1$ when densities are biomass or energy & M & Predation changes $x+y$ by $-(1-e)fy$. & \checkmark\,S13a & -- \\
N2 & $m>0$ & D & Natural death rate. & -- & -- \\
N3 & $\psi>0$; predator Allee makes $y=0$ locally stable, unless predator immigration (a small stable equilibrium instead) or constant-yield harvest ($y=0$ reached in finite time); $T_3=e\,q(y)\,T_2$ & M & Mate-finding factor $q=y/(y+\psi)$ \cite{courchamp2008}. & \checkmark\,S16, S32b & Marked as an exception to Rule~1 ($\S$). Second evaluation: the note claimed local stability of $y=0$ also with predator immigration. \\
N4 & $\tau\ge 0$; initial data on $[-\tau,0]$; $T_3(t)=e\,s\,T_2(t-\tau)$ & D & \cite{cushing1977book,kuang1993}. & -- & Marked as an exception to Rule~1 ($\ddagger$). \\
N5 & Delay-dependent parameters: stability switches, often stable, unstable, stable again; oscillations in a window of maturation delays & Q & \cite{beretta2002,gourley2004}. & \cite{beretta2002}: abstract; \cite{gourley2004}: abstract & Source added. \\
N6 & $\delta^{J}\ge 0$; juvenile survival $\exp(-\delta^{J}\tau)$ & M & Solution of $J'=-\delta^{J}J$ over the delay \cite{gourley2004}. & \checkmark\,S15 & -- \\
N7 & $\delta\ge 0$ & D & Density-dependent predator mortality \cite{bazykin1998}. & -- & -- \\
N8 & $\beta_{jj'}\ge 0$ & D & Interference competition between predators. & -- & -- \\
N9 & Leslie--Gower: $s,\eta>0$; $\eta$ = predators supported per unit prey & M & Predator equilibrium $y=\eta x$ \cite{leslie1960}. & \checkmark\,S18c & -- \\
N10 & Leslie--Gower singular without prey; $C>0$ removes it & M & \cite{azizalaoui2003}. & \checkmark\,S18a, S18b & -- \\
\addlinespace
\multicolumn{6}{@{}l}{\textit{Predators eating predators}}\\[2pt]
Y1 & $a^{y},h^{y},\rho^{y},\lambda^{y},\eta^{y}>0$ & D & As for prey links, indexed by eaten and eating predator. & -- & -- \\
Y2 & $e^{y}>0$; $e^{y}<1$ when densities are biomass or energy & M & Mass balance as in N1. & \checkmark\,S13a & -- \\
Y3 & Cannibalism: $e^{y}_{ll}\le 1$ in any units; if violated, blow-up in finite time for $h=0$, at most exponential growth for $h>0$ & M & Unstructured model: numbers proportional to biomass; net link $(e^{y}-1)a^{y}y^{2}/D$. & \checkmark\,S13a, S14a, S14b, S14c & The blow-up was stated without the restriction $h=0$; with handling time it cannot occur (S14c). \\
Y4 & Cannibalism can stabilise and raise both standing stocks (one prey, one predator) & Q & \cite{kohlmeier1995}. & \cite{kohlmeier1995}: abstract & Second finding added; restricted to one prey and one predator at the second evaluation. \\
Y5 & Intraguild predation: coexistence needs the intraguild prey to be the better competitor and the intraguild predator to gain substantially; most likely at intermediate productivity & Q & \cite{holt1997}; review \cite{polis1989}. & \cite{holt1997}: abstract; \cite{polis1989}: not accessed & ``Only at intermediate productivity\textquotedblright{} and the missing second condition corrected. \\
Y6 & Holling II food chains can be chaotic & Q & \cite{hastings1991}. & \cite{hastings1991}: abstract & -- \\
Y7 & Reciprocal predation between two predators: existence and stability of equilibria, Hopf bifurcation at the two positive equilibria & Q & \cite{almohanna2025}. & \cite{almohanna2025}: abstract & -- \\
Y8 & No cannibalism with Leslie--Gower & C & Builder choice: the predator's capacity would depend on its own density. & -- & -- \\
Y9 & Refuge, habitat, cooperation, herd and group defence act on prey links only & C & Builder choice, stated in a note. & -- & -- \\
\addlinespace
\multicolumn{6}{@{}l}{\textit{Added terms ($\Phi$)}}\\[2pt]
H1 & $q,E,H\ge 0$; $\ell$ constants $>0$ & D & \cite{clark1976,brauer1979,gupta2013,chaudhuri1996}. & -- & -- \\
H2 & Effort harvest: prey alone persists iff $qE<r$ (without prey immigration) & M & Exponential, logistic, $\theta$-logistic and delayed logistic prey; equilibrium $K(1-qE/r)^{1/\theta}$. & \checkmark\,S19a, S19b & Second evaluation: not stated with immigration, which keeps the prey alive at any effort. \\
H3 & Constant yield, logistic prey: a positive equilibrium needs $H\le rK/4$ and a stable one $H<rK/4$ (without prey immigration or seasonal forcing) & M & Maximum of $rx(1-x/K)$ is $rK/4$ at $K/2$; the larger equilibrium is stable \cite{brauer1979}. & \checkmark\,S20, S20b & Second evaluation: not stated with immigration or seasonal forcing, which change the threshold. \\
H4 & Constant yield must stop at zero density & M & At $x=0$ the right side is $-H<0$. & \checkmark\,S21 & -- \\
H5 & $I\ge 0$; immigration makes the Lotka--Volterra centre asymptotically stable (stated only for the classical model) & M, A & Trace $=-I/x^{*}<0$ (derived). Tahara et al.: analytic for linear predation, under stated conditions for types II and III. & \checkmark\,S22; \cite{tahara2018}: full text (Eq.~(1), Case A1 ($\gamma=bc/2m$); Eq.~(4), Table~1 and Suppl.\ Tables S2--S5) & ``For three types of functional response\textquotedblright{} (abstract) qualified after reading the full text. Second evaluation: shown only when the rest of the model is classical. \\
H6 & $\zeta\ge 0$ (prey kill predators without gain) & D & \cite{tang2015}. & -- & -- \\
H7 & Prey eat predators: $e^{x}>0$; $e^{x}<1$ in biomass units & M & $x+y$ changes by $-(1-e)axy-(1-e^{x})\zeta xy$. & \checkmark\,S13b & -- \\
H8 & $\vartheta\ge 0$, $g>0$ & D & Toxic prey \cite{chattopadhyay2002}. & -- & -- \\
H9 & Predator harvesting ($\Phi_y$.1) with logistic prey and linear predation: the predator persists iff $m+q_2E_2<eaK$ & M & Interior equilibrium $x^{*}=(m+q_2E_2)/(ea)$, $y^{*}=r(1-x^{*}/K)/a>0$ iff $x^{*}<K$; model of \cite{chaudhuri1996}. & \checkmark\,S42a & Added 19 October 2026 (audit of the entries table). \\
H10 & Anti-predator attack ($\Phi_y$.2): the predator cannot persist if $\zeta\ge ea$ & M & Per-capita predator growth $(ea-\zeta)x-m\le -m<0$ for every $x\ge 0$; model of \cite{tang2015}. & \checkmark\,S42b & Added 19 October 2026 (audit of the entries table). \\
\addlinespace
\multicolumn{6}{@{}l}{\textit{New state variables ($\mathbf z$)}}\\[2pt]
Z1 & Competing predators (z.9): with one prey, only the predator with the lowest break-even prey density ($m_i/(e_ia_i)$ for linear responses) persists at equilibrium & M, Q & At the prey level $x^{*}=m_1/(e_1a_1)$ where predator 1 is at equilibrium, predator 2 has per-capita growth $e_2a_2x^{*}-m_2<0$ iff $m_1/(e_1a_1)<m_2/(e_2a_2)$ (derived); the exclusion principle for saturating responses is \cite{hsu1978}, and \cite{armstrong1980} give the conditions for coexistence on one prey. & \checkmark\,S42d; \cite{hsu1978}: abstract; \cite{armstrong1980}: abstract & Added 19 October 2026 (audit of the entries table). \\
Z2 & Metapopulation (z.16): the occupied fraction persists iff $c>\epsilon$ & M & Equilibrium $p^{*}=1-\epsilon/c$ is positive iff $c>\epsilon$, and then stable ($\mathrm d\dot p/\mathrm dp=\epsilon-c<0$); $p=0$ is stable iff $c<\epsilon$; model of \cite{levins1969}. & \checkmark\,S42e & Added 19 October 2026 (audit of the entries table). \\
\addlinespace
\multicolumn{6}{@{}l}{\textit{Framework and general statements}}\\[2pt]
F1 & $D^{x},D^{y}>0$; no-flux boundary conditions & D & \cite{skellam1951}. & -- & -- \\
F2 & Turing instability (one prey, one predator, no delay or seasonal forcing) iff $D^{x}D^{y}\mu_k^{2}-(D^{y}J_{11}+D^{x}J_{22})\mu_k+\det J<0$ for a no-flux eigenvalue $\mu_k>0$ of $-\Delta$ on $\Omega$; such a band exists iff $D^{y}J_{11}+D^{x}J_{22}>0$ and $(D^{y}J_{11}+D^{x}J_{22})^{2}>4D^{x}D^{y}\det J$; impossible for $J_{22}=0$; possible with $J_{22}<0$ (faster predator) or $J_{22}>0$ (faster prey) & M & Dispersion relation with $k^{2}=\mu_k$; self-limitation \cite{segel1972}; cooperation example. More species: an eigenvalue of $J-\mu_kD$ with positive real part. & \checkmark\,S27a, S27b, S27c, S27d, S27e, S27f; \cite{segel1972}: not accessed & ``Turing patterns need predator self-limitation\textquotedblright{} was wrong: predator self-enhancement with faster prey diffusion also works (S27f). Second evaluation: the bounded-domain eigenvalue condition added; the two-species conditions are no longer stated for food webs, delays or seasonal forcing. \\
F3 & It\^o noise ($\sigma>0$), prey alone: stationary density iff $r>\sigma^{2}/2$ (logistic, $\theta$-logistic), $r-d>\sigma^{2}/2$ (fear, $c^{x}>0$), always (Gompertz); extinction a.s.\ when reversed; not stated with prey immigration or seasonal forcing. Positivity: not with constant yield, and the herd loss $x^{p}$ is not covered & M & Zero-flux stationary densities; comparison for $d\ln x$; noise as in \cite{rudnicki2003}. & \checkmark\,S26a, S26b, S26c, S26d, S26e, S26f, S26g; \cite{rudnicki2003}: not accessed & ``Persists\textquotedblright{} was undefined and rested on the low-density drift only; now defined by the stationary density. Second evaluation: exceptions to positivity and to the prey-alone results added. \\
F4 & Non-negative orthant invariant, except with constant-yield harvest & M & Quasi-positivity of every term except $-H$. & \checkmark\,S30, S21 & -- \\
F5 & Parameters non-negative unless their constraint says otherwise & D & The weak-Allee threshold $A$ may be negative. & -- & General note corrected. \\
F6 & Regularity: Gompertz unique; herd unique forward only; interference not unique forward; ratio dependence Lipschitz; Leslie--Gower singular & M & Table~\ref{M-tab:regularity}. & \checkmark\,S7b, S8d, S11d, S28b, S18a & Rule~4 listed three cases under one heading; now classified. \\
F7 & Demographic noise of order $N^{-1/2}$ sustains large quasi-cycles around a stable equilibrium & Q & Individual-level model; Volterra mean field with a stable equilibrium; resonant power spectrum \cite{mckane2005}. & \cite{mckane2005}: full text (arXiv:\allowbreak q-bio/\allowbreak 0501023, Eqs.~(3), (6), (9); Fig.~2) & -- \\
\addlinespace
\multicolumn{6}{@{}l}{\textit{Classical results}}\\[2pt]
K1 & Lotka--Volterra: centre at $(m/ea,\,r/a)$, first integral $H$ & M & \cite{lotka1925,volterra1926}. & \checkmark\,S23 & -- \\
K2 & Logistic prey, Holling I: interior equilibrium iff $m<eaK$, globally stable & M & $\dot V=-(er/K)(x-x^{*})^{2}\le 0$ and LaSalle. & \checkmark\,S24 & -- \\
K3 & Rosenzweig--MacArthur: $x^{*}=m/[a(e-mh)]$, exists iff $e>mh$, $x^{*}<K$; then hyperbolic and stable iff $x^{*}>\tfrac12(K-1/ah)$, unstable below, weakly attracting at equality (K4) & M & Trace $\propto(x_{\rm hump}-x^{*})$; $\det J=mr[Ka(e-mh)-m]/(Kae)$ \cite{rosenzweig1963}. & \checkmark\,S25a, S25b, S25c & The sign of $\det J$ was checked at random parameter values; now proved in general. \\
K4 & Supercritical Hopf at $x^{*}=x_{\rm hump}$; unique limit cycle beyond it; paradox of enrichment & M, Q & $16\omega^{2}\ell_1=-4\beta s/[(1+s)^{2}(1+2s)^{2}]$; cross-check with a second formula. Limit cycle \cite{cheng1981}; enrichment \cite{rosenzweig1971}. & \checkmark\,S25d0, S25d, S25e; \cite{cheng1981}: abstract; \cite{rosenzweig1971}: abstract & The sign of $\ell_1$ was sampled at 12 parameter sets; now proved in general. \\
\addlinespace
\multicolumn{6}{@{}l}{\textit{Entries added after the scoping review}}\\[2pt]
X1 & Food-limited growth (T1.12): $c>0$; $c\to\infty$ gives the logistic law; the per-capita rate falls from $r$ at $x=0$ to $0$ at $x=K$, so the prey-only equilibrium stays $K$ & D, M & Replacement rate at saturation $c$; form of \cite{smith1963}, Eq.~(5); values, limit and derivative of the per-capita rate. & \checkmark\,S34a, S34b, S34c; \cite{smith1963}: full text (Eq.~(5) and definitions, pp.~652--653) & Added 19 October 2026 (scoping review, Stage C). \\
X2 & Sexual reproduction (T3.7): births in proportion to encounters of males and females, $cy^{2}$, with a Leslie--Gower loss $wy^{2}/(x+D)$ & A & Top predator of the food chain in \cite{upadhyay1997}, written here for one prey and one predator. & \cite{upadhyay1997}: full text (Eq.~(1c) and text, p.~1934) & Added 19 October 2026. \\
X3 & Sexual reproduction (T3.7): the per-capita gain is positive iff $x>w/c-D$; if $x\ge x_0>w/c-D$ for all $t$, $y$ blows up in finite time & M & Comparison with $\dot y=ky^{2}$, $k=c-w/(x_0+D)>0$. & \checkmark\,S35a, S35b & Added 19 October 2026. \\
X4 & Allelopathy ($\Phi_x$.7): a predator toxin kills prey at per-capita rate $\gamma xy^{2}$; between two competitors the per-capita rates are $\gamma_iN_1N_2$ & Q & \cite{fitria2015,chattopadhyay1996}. & \cite{fitria2015}: full text (Eq.~(1), p.~59); \cite{chattopadhyay1996}: full text (Eq.~(1), p.~287) & Added 19 October 2026. \\
X5 & Advection (S.7): logistic prey alone spread upstream iff $v<2\sqrt{rd_1}$ & A & Upstream speed $c^{*}-v$ with the Fisher speed $c^{*}=2\sqrt{rd_1}$ \cite{hilker2010}. & \cite{hilker2010}: full text (Eqs.~(1)--(2) and invasion speeds, pp.~176--177) & Added 19 October 2026. \\
X6 & Spatial heterogeneity (S.8): with lethal boundaries, logistic prey with growth rate $r(\xi)$ of either sign persist if $1/d_1>\lambda_1^{+}(r)$; fragmenting favourable habitat hinders persistence & A & Stated in \cite{cantrell1989} for $\Delta u+\lambda u(m-u)=0$; here $\lambda=1/d_1$. Predator--prey version in \cite{du2004}. & \cite{cantrell1989}: abstract; \cite{du2004}: full text (model (1.1), p.~332) & Added 19 October 2026. \\
X7 & Feedback control (M.3), logistic prey: the equilibrium falls to $K/(1+cqK/f)$ and is locally asymptotically stable & M & Model of \cite{gopalsamy1993}; Jacobian at the equilibrium. & \checkmark\,S36a, S36b; \cite{gopalsamy1993}: full text (model in the abstract and Sec.~1) & Added 19 October 2026. \\
X8 & Hybrid control (M.3): $\alpha F+(1-\alpha)\,\mathrm{id}$ keeps the fixed points and maps each eigenvalue $\lambda$ to $\alpha\lambda+1-\alpha$ & M & Form of \cite{luo2003}. & \checkmark\,S36c; \cite{luo2003}: full text (Eq.~(2), p.~776) & Added 19 October 2026. \\
X9 & Z-type control (M.3): the deviation $x-x_d$ from a target decays exponentially & Q & \cite{lacitignola2016}. & \cite{lacitignola2016}: full text (Sec.~2.1, Eq.~(2), p.~11) & Added 19 October 2026. \\
X10 & Economic profit constraint (M.4): with one harvested species and zero profit, the effort settles where rent vanishes, $x=c_x/p_x$ & M & Bionomic equilibrium of \cite{gordon1954}; algebraic equation of \cite{zhang2009bioecon}. & \checkmark\,S37; \cite{zhang2009bioecon}: full text (model (1), p.~579) & Added 19 October 2026. \\
X11 & Economic profit constraint (M.4): a singularity-induced bifurcation occurs at a critical profit & Q & \cite{zhang2009bioecon}. & \cite{zhang2009bioecon}: full text (abstract and Sec.~3) & Added 19 October 2026. \\
X12 & Mutualist of the prey (z.20): predation decreases with the mutualist, which gains from the prey and is self-limited; its effect on prey growth may have either sign & Q & \cite{rai1983}. & \cite{rai1983}: full text (hypotheses (H1), (G1), (P1), pp.~16--17) & Added 19 October 2026. \\
X13 & A protective population that lowers predation on the prey (human shield) is a special case of z.20 & C & Our classification in the scoping review (Stage C). & -- & Added 19 October 2026. \\
X14 & Scavenger (z.21): gains from kills and natural deaths at no cost to $x$ or $y$; self-limitation keeps it bounded; it may also prey on $x$ & Q & \cite{nolting2008,previte2013}. & \cite{nolting2008}: full text (Eq.~(3), p.~217); \cite{previte2013}: full text (Eq.~(3.1), p.~528) & Added 19 October 2026. \\
X15 & Sex structure (z.22): one sex produces both; the sexes share density dependence and are eaten alike & Q & Model of Liu (2005), quoted in \cite{xiong2007}. & \cite{xiong2007}: full text (model (1.1), p.~1214) & Added 19 October 2026. \\
X16 & Lateral dimorphism (z.23): with rates symmetric between morphs, the symmetric state is invariant & M & Model 1 of \cite{nakajima2004}. & \checkmark\,S38; \cite{nakajima2004}: full text (Model~1, p.~693) & Added 19 October 2026. \\
X17 & Lateral dimorphism (z.23): when cross predation dominates, the morph frequencies oscillate around $1/2$, and a missing morph can invade & Q & \cite{takahashi1994,nakajima2004}. & \cite{nakajima2004}: full text (abstract, p.~692); \cite{takahashi1994}: abstract & Added 19 October 2026. \\
X18 & Power-law interactions (T2.16): Gilpin--Ayala predator--prey systems raise each density in each per-capita rate to its own power (prey loss $ay^{\beta}$, predator gain $e\,x^{p}$, self-limitation $\delta y^{q}$) & Q & Discrete system of \cite{liao2007gilpin}; the continuous system of \cite{chen2005gilpin} as quoted there. & \cite{liao2007gilpin}: full text (systems (1.1) and (1.6), pp.~500--501) & Added 19 October 2026 (adjudication of the blind check). \\
X19 & Power-law interactions (T2.16): with $0<p<1$ the predator gain exceeds the prey loss at low prey density, so in biomass units the model creates biomass & M & Ratio gain/loss $=(e/a)x^{p-1}y^{1-\beta}\to\infty$ as $x\to0^{+}$ (Rule~2). & \checkmark\,S39 & Added 19 October 2026. \\
X20 & Independent numerical response (T3.8): the predator's demographic rates can be modelled as their own function of food availability (demographic numerical response) & Q & \cite{bayliss2002}. & \cite{bayliss2002}: abstract & Added 19 October 2026 (adjudication of the blind check). \\
X21 & Independent numerical response (T3.8): in biomass units the gain may not exceed the intake, $e\,g\le f$; a Monod--Haldane gain with a Beddington--DeAngelis loss violates this as $y\to\infty$ & M & Rule~2; limit of $g/f$ as $y\to\infty$. & \checkmark\,S40 & Added 19 October 2026. \\
X22 & Starvation mortality (T4.4): extra death rate $T^{-1}\ln(\bar f/f)$ of a predator receiving the ration $f\le\bar f$; zero at satiation, decreasing in the ration, unbounded as $f\to0$ & Q, M & Form of \cite{arditi1978}, adapted there from an ecosystem model; properties derived. & \checkmark\,S41a; \cite{arditi1978}: full text (Eqs.~(10)--(13), pp.~179--180) & Added 19 October 2026 (adjudication of the blind check). \\
X23 & Starvation mortality (T4.4): at a constant ration $f_1$ and without births, the hunger term alone shrinks the predator by the factor $f_1/\bar f$ in time $T$ & Q, M & Stated by \cite{arditi1978} as the reason for the form; verified. & \checkmark\,S41b; \cite{arditi1978}: full text (Eq.~(14) and Fig.~3, pp.~180--181) & Added 19 October 2026. \\
X24 & Self- and cross-diffusion (S.2): self-diffusion $\alpha\ge0$ and cross-diffusion $\beta\ge0$ & D & Broadened from cross-diffusion alone when the author adjudicated the blind check; model of \cite{shigesada1979}. & -- & S.2 broadened 19 October 2026. \\
X25 & Taxis (S.3): the prey-taxis term also models immature prey moving toward higher densities of mature prey & Q & \cite{li2023taxis}. & \cite{li2023taxis}: full text (model (3) and its explanation, p.~18593) & S.3 broadened 19 October 2026. \\
\addlinespace
\end{longtable}
\endgroup
